\documentclass[conference]{IEEEtran}
\IEEEoverridecommandlockouts

\usepackage{cite}
\usepackage{amsmath,amssymb,amsfonts}
\usepackage{booktabs}
\usepackage{url}

\def\BibTeX{{\rm B\kern-.05em{\sc i\kern-.025em b}\kern-.08em
    T\kern-.1667em\lower.7ex\hbox{E}\kern-.125emX}}

\begin{document}

\title{Theoretical Foundations for Speed and Torque Control of a Switched Reluctance Motor}

\author{%
\IEEEauthorblockN{Ahmed Mostafa}
\IEEEauthorblockA{\textit{Dept. of Mechatronics Engineering}\\
\textit{German University in Cairo}\\
Cairo, Egypt\\
Student ID: 55-1591}
\and
\IEEEauthorblockN{Andrew Abdelmalak}
\IEEEauthorblockA{\textit{Dept. of Mechatronics Engineering}\\
\textit{German University in Cairo}\\
Cairo, Egypt\\
Student ID: 55-22771}
\and
\IEEEauthorblockN{Adham Bassem}
\IEEEauthorblockA{\textit{Dept. of Mechatronics Engineering}\\
\textit{German University in Cairo}\\
Cairo, Egypt\\
Student ID: 55-21599}
\and
\IEEEauthorblockN{Ahmed Mansour}
\IEEEauthorblockA{\textit{Dept. of Mechatronics Engineering}\\
\textit{German University in Cairo}\\
Cairo, Egypt\\
Student ID: 55-0253}
}

\maketitle

\begin{abstract}
Switched reluctance motors (SRMs) are attractive for electric-drive applications due to robust
construction, absence of rotor magnets, and wide speed capability. Their nonlinear inductance
profile and commutation behavior, however, make high-performance torque and speed control
challenging. This paper consolidates the theoretical background for an SRM speed/torque control
project milestone. We summarize machine fundamentals, nonlinear dynamic equations,
converter operation, and established control methods including hysteresis current control,
fixed-frequency PWM, turn-on/turn-off angle scheduling, and torque sharing functions (TSF).
A cascaded control architecture is presented as a practical implementation template for
subsequent MATLAB/Simulink realization.
\end{abstract}

\begin{IEEEkeywords}
Switched reluctance motor, electric drives, torque ripple, torque sharing function, PI speed control
\end{IEEEkeywords}

\section{Introduction}
SRMs are increasingly considered for applications that prioritize reliability, fault tolerance,
and reduced dependence on rare-earth materials~\cite{Miller1993,Krishnan2001}. The machine
structure is simple---a salient-pole stator carrying concentrated windings and a salient-pole
rotor with neither magnets nor windings---but control complexity is high because torque is
generated through position-dependent inductance rather than sinusoidal back-EMF. As a result,
both commutation logic and current regulation must be tightly synchronized with rotor
position~\cite{Husain1996}.

Early analytical treatments demonstrated that the SRM's nonlinearity is amenable to
feedback-linearizing control~\cite{Ilic1987}, while subsequent work focused on practical
converter topologies, current regulation, and torque-ripple suppression. This milestone
addresses the theoretical layer that precedes implementation. The objective is to establish a
coherent model-to-control chain---from machine equations through converter switching to torque
sharing and speed regulation---that can be translated into MATLAB/Simulink simulation and,
subsequently, real-time drive experiments.

\section{SRM Construction and Inductance Profile}
A common SRM configuration is the three-phase 6/4 geometry: six stator poles each carrying a
concentrated coil, and four rotor poles with no windings or magnets. Diametrically opposite
stator coils are connected in series to form one phase, yielding three independent phases
($A$, $B$, $C$) displaced by $30^{\circ}$ electrical~\cite{Miller1993}.

The phase inductance $L(\theta)$ varies periodically with rotor position between two extremes.
At the \emph{aligned} position ($\theta_a$), a rotor pole is fully aligned with an energized stator pole
pair, the magnetic reluctance of the airgap path is minimized, and inductance reaches its
maximum $L_{\max}$. At the \emph{unaligned} position ($\theta_u$), the rotor interpolar axis faces the
stator pole, reluctance is maximized, and inductance has its minimum $L_{\min}$.

Positive motoring torque is produced only in the rising-inductance interval ($dL/d\theta > 0$),
hence each phase must be energized during that region and demagnetized before the
falling-inductance zone. The ratio $L_{\max}/L_{\min}$ directly influences torque capability;
values of 6--10 are typical for well-designed machines~\cite{Krishnan2001}.

Under magnetic saturation the inductance becomes a function of both $\theta$ and $i$, i.e.\
$L(\theta, i)$. A nonlinear look-up table, typically obtained from finite-element analysis (FEA),
is therefore required for accurate simulation.

\section{SRM Dynamic Model}
\subsection{Electrical and Magnetic Relations}
For phase voltage $v_k$, current $i_k$, resistance $R$, and flux linkage $\psi_k$, the standard
phase relation is
\begin{equation}
v_k = R i_k + \frac{d\psi_k}{dt}.
\end{equation}
With position-dependent inductance, this expands to
\begin{equation}
v_k = R i_k + L(\theta, i_k)\frac{di_k}{dt} + i_k\frac{\partial L}{\partial \theta}\omega,
\label{eq:phase_voltage}
\end{equation}
where $\omega$ is rotor speed. The last term represents motional EMF and grows with speed,
eventually preventing the controller from maintaining the desired current waveform above a
characteristic \emph{base speed}.

Under linear (unsaturated) assumptions, flux linkage reduces to
$\psi_k(\theta, i_k) = L(\theta)\,i_k$.
For practical SRMs, saturation requires the nonlinear map $\psi(\theta, i)$, commonly implemented
as two-dimensional lookup tables derived from FEA or measurements~\cite{Krishnan2001}.

\subsection{Torque and Mechanical Dynamics}
The phase torque expression under the linear inductance assumption is
\begin{equation}
T_{e,k} = \frac{1}{2}\, i_k^2 \frac{dL_k(\theta)}{d\theta},
\end{equation}
and total electromagnetic torque is $T_e = \sum_k T_{e,k}$. In the saturated regime, torque is
more accurately obtained from the co-energy integral. Mechanical motion obeys
\begin{equation}
J\frac{d\omega}{dt} = T_e - T_L - B\omega,
\label{eq:mechanical}
\end{equation}
where $J$ is rotor inertia, $T_L$ the load torque, and $B$ a viscous damping coefficient.

\section{Converter Topology}
The asymmetric half-bridge converter remains the dominant SRM topology because each phase can
be energized and demagnetized independently using two switches and two
diodes~\cite{Miller1993}. Three switching states are available per phase, summarized in
Table~\ref{tab:converter}.

\begin{table}[htbp]
\caption{Asymmetric Half-Bridge Switching States}
\label{tab:converter}
\centering
\begin{tabular}{@{}lccc@{}}
\toprule
Mode & Q1 & Q2 & Phase voltage \\
\midrule
Magnetisation  & ON  & ON  & $+V_{dc}$ \\
Freewheeling   & ON  & OFF & $\approx 0$ \\
Demagnetisation & OFF & OFF & $-V_{dc}$ \\
\bottomrule
\end{tabular}
\end{table}

Fast current decay during demagnetisation is essential: applying the full negative DC-link
voltage forces the current to zero quickly, preventing negative torque in the
falling-inductance region. The speed at which demagnetisation completes limits the motor's
maximum operating speed and determines the required turn-off angle advance~\cite{Husain1996}.

\section{Current Control Strategies}
\subsection{Hysteresis Current Control}
At low-to-medium speed the back-EMF is small and current must be actively limited. Hysteresis
(bang-bang) control compares the reference current $i^*$ with the measured phase current using a
hysteresis band of width $\Delta i$:
\begin{itemize}
\item $i < i^* - \Delta i/2$: magnetisation (both switches ON).
\item $i > i^* + \Delta i/2$: demagnetisation (both switches OFF).
\item Otherwise: freewheeling (maintain previous state).
\end{itemize}
The switching frequency is variable, depending on DC-link voltage, inductance, and band width.
A narrower band yields better current tracking but increases switching losses; 5--10\% of rated
current is a common design choice~\cite{Husain1996}.

\subsection{PWM Current Control}
An alternative is fixed-frequency PWM applied to one switch while the other remains closed.
The average phase voltage is $V_{\mathrm{avg}} = D\cdot V_{dc}$, where $D\in[0,1]$ is the duty
cycle modulated by a PI current controller. PWM control offers a predictable switching
frequency, simplifying EMI-filter design, but requires a current-feedback loop for effective
regulation~\cite{Krishnan2001}.

\subsection{Turn-On / Turn-Off Angle Control}
Above the base speed, the back-EMF prevents hysteresis or PWM controllers from maintaining the
desired current, and the motor operates in \emph{single-pulse} (voltage-control) mode. The
primary control variables become:
\begin{itemize}
\item \textbf{Turn-on angle} $\theta_{\mathrm{on}}$: advancing it allows more time for current to build.
\item \textbf{Turn-off angle} $\theta_{\mathrm{off}}$: advancing it reduces negative torque; delaying it
      increases average torque but risks torque reversal.
\end{itemize}
Optimal scheduling of $(\theta_{\mathrm{on}}, \theta_{\mathrm{off}})$ as a function of speed and
load can be pre-computed and stored in a two-dimensional look-up table~\cite{Miller1993}.

\section{Torque Ripple and Torque Sharing Functions}
\subsection{Sources of Ripple}
Torque ripple is a principal SRM drawback, arising from: (i)~discrete commutation between
phases, creating torque dips; (ii)~magnetic saturation limiting torque near the aligned
position; and (iii)~the inherently nonuniform $dL/d\theta$ profile. Reducing ripple is critical
for applications requiring smooth speed regulation or low acoustic
noise~\cite{Schramm1992,Xue2009}.

\subsection{TSF Formulations}
TSF methods distribute the total torque reference $T^*$ across overlapping phases so that their
sum remains constant during commutation:
\begin{equation}
T_k^*(\theta) = \bigl(1-f_r(\theta)\bigr)T^*,\quad T_{k+1}^*(\theta)=f_r(\theta)\,T^*,
\end{equation}
where $f_r(\theta)$ is the rising sharing function defined over the overlap angle
$[\theta_{\mathrm{on}},\;\theta_{\mathrm{on}}+\theta_{\mathrm{ov}}]$. Table~\ref{tab:tsf}
compares common formulations.

\begin{table}[htbp]
\caption{Common Torque Sharing Function Formulations}
\label{tab:tsf}
\centering
\begin{tabular}{@{}lll@{}}
\toprule
TSF type & $f_r(\theta)$ & Key trade-off \\
\midrule
Linear      & $(\theta-\theta_{\mathrm{on}})/\theta_{\mathrm{ov}}$ & Simple; large $di/d\theta$ at edges \\
Sinusoidal  & $\tfrac{1}{2}(1-\cos\tfrac{\pi\theta}{\theta_{\mathrm{ov}}})$ & Smooth; moderate current stress \\
Cubic       & $3x^2-2x^3,\; x=\theta/\theta_{\mathrm{ov}}$ & Zero slope at boundaries \\
Exponential & $1-e^{-\alpha\theta}$ & Tracks saturation; needs tuning \\
\bottomrule
\end{tabular}
\end{table}

The feasibility of a TSF is ultimately limited by the available DC-link voltage: if the demanded
$di/dt$ exceeds $V_{dc}/L(\theta)$, actual torque deviates from reference and ripple
increases~\cite{Xue2009,Zhu2015}. At high speeds, even smooth TSFs become infeasible, and
advanced methods such as optimal current profiling~\cite{Schramm1992} or model-predictive
control are employed. Individual phase torque references are inverted through a nonlinear
$T(\theta, i)$ look-up table to obtain corresponding current references $i_k^*$, tracked by the
inner current loop.

\section{Cascaded Speed Control Architecture}
A standard outer-loop PI controller generates the total torque demand:
\begin{equation}
T^*(s)=K_p\,e_\omega(s)+\frac{K_i}{s}\,e_\omega(s),
\end{equation}
with speed error $e_\omega = \omega^* - \omega$. An anti-windup mechanism is essential because
the torque reference is saturated at a physical limit $T_{\max}$; without it, integrator
windup causes significant overshoot during transient loading.

The TSF block distributes $T^*$ among phases, the torque-to-current inversion obtains phase
current references, and the inner current loop (hysteresis or PWM) tracks them via the
converter. Position feedback---from an encoder or sensorless observer---is required at every
level: the commutation logic, TSF indexing, and torque-to-current inversion all depend on
$\theta$~\cite{Cheok2000}. This tight coupling between position knowledge and control
performance underscores why accurate and fast position sensing is indispensable for SRM drives.

Gain tuning is typically performed via simulation-based trial-and-error or linearised
small-signal models. A bandwidth separation of at least one decade between the inner current
loop and the outer speed loop is recommended to preserve stability~\cite{Ilic1987,Krishnan2001}.

\section{Discussion and Next Milestone}
The available project evidence is documentation-centric and does not yet include \texttt{.m}
scripts, Simulink models, or waveform outputs. This paper therefore limits claims to
theoretical readiness. The next milestone should provide: (1)~a nonlinear SRM model
implementation in MATLAB/Simulink, (2)~baseline hysteresis and PWM controllers, (3)~a TSF
comparison under common operating points, and (4)~closed-loop speed transient metrics
including rise time, overshoot, and steady-state error.

\section{Conclusion}
A complete theory baseline for SRM speed and torque control has been consolidated from project
material and canonical references. The paper has detailed the machine construction and
inductance characteristics, derived the governing electrical and mechanical equations,
catalogued converter switching states, surveyed three current-control strategies (hysteresis,
PWM, angle control), and presented TSF-based torque-ripple mitigation together with a cascaded
PI speed-control architecture. These results provide a rigorous and traceable starting point for
model-based implementation in subsequent milestones.

\bibliographystyle{IEEEtran}
\bibliography{references}

\end{document}
